\Needspace{6\baselineskip}
\begin{center}\hypertarget{native-259-article186}{}\pdfbookmark[2]{186.}{article-186}186.\end{center}
Im Übrigen lässt sich der Grundsatz, dass die Quadrate der Differenzen zwischen den beobachteten und den berechneten Größen eine möglichst kleine Summe ergeben sollen, auch unabhängig von der Wahrscheinlichkeitsrechnung in folgender Weise betrachten.

Ist die Zahl der Unbekannten gleich der Zahl der voneinander unabhängigen beobachteten Größen, so lassen sich die Unbekannten so bestimmen, dass sie den Beobachtungen genau genügen. Ist die Zahl der Unbekannten dagegen kleiner, so ist eine vollkommen genaue Übereinstimmung insoweit unerreichbar, als die Beobachtungen keine absolute Präzision haben. In diesem Fall müssen wir also die bestmögliche Übereinstimmung anstreben oder, anders gesagt, die Differenzen möglichst klein machen. Dieser Begriff ist jedoch von Natur aus etwas unbestimmt. Ein System von Werten für die Unbekannten, das \emph{alle} Differenzen jeweils kleiner macht als ein anderes, verdient zweifellos den Vorzug. Gleichwohl bleibt die Wahl zwischen zwei Systemen, von denen das eine für manche Beobachtungen eine bessere Übereinstimmung liefert und das andere für andere, in gewissem Maße unserem Ermessen überlassen. Offenbar lassen sich unzählige verschiedene Grundsätze angeben, durch die die erste Bedingung erfüllt wird. Bezeichnen wir die Differenzen zwischen Beobachtung und Rechnung mit $\Delta$, $\Delta'$, $\Delta''$ usw. Die erste Bedingung ist sowohl erfüllt, wenn $\Delta\Delta+\Delta'\Delta'+\Delta''\Delta''+\text{usw.}$ minimal wird (darin besteht unser Grundsatz), als auch, wenn $\Delta^4+\Delta'^4+\Delta''^4+\text{usw.}$ oder $\Delta^6+\Delta'^6+\Delta''^6+\text{usw.}$ oder allgemein die Summe der Potenzen mit einem beliebigen geraden Exponenten ein Minimum annimmt. Von all diesen Grundsätzen ist jedoch unserer der einfachste, da die übrigen zu äußerst verwickelten Rechnungen führen würden. Unseren Grundsatz, den wir bereits seit 1795 anwenden, hat kürzlich auch der angesehene \textsc{Legendre} in seinem Werk \emph{Nouvelles méthodes pour la
\SourceLeaf{260}détermination des orbites des comètes, Paris} 1806 ausgesprochen. Dort sind mehrere weitere Eigenschaften dieses Grundsatzes dargelegt, die wir hier der Kürze halber übergehen.

Wählten wir eine Potenz mit einem unendlich großen geraden Exponenten, so gelangten wir zu dem System, in dem die größten Differenzen möglichst klein werden.

Zur Lösung linearer Gleichungen, deren Zahl größer ist als die Zahl der Unbekannten, verwendet der berühmte \textsc{Laplace} einen anderen Grundsatz, den schon der angesehene \textsc{Boscovich} vorgeschlagen hatte. Die Differenzen selbst, alle mit positivem Vorzeichen genommen, sollen die kleinste Summe ergeben. Es lässt sich leicht zeigen, dass ein allein nach diesem Grundsatz ermitteltes System von Werten für die Unbekannten notwendigerweise\footnote{Ausgenommen sind Sonderfälle, in denen die Lösung in gewissem Maße unbestimmt bleibt.} ebenso vielen der gegebenen Gleichungen exakt genügen muss, wie es Unbekannte gibt. Die übrigen Gleichungen kommen also nur insoweit in Betracht, als sie \emph{die Wahl mitbestimmen}. Gehört etwa die Gleichung $V=M$ zu denen, die nicht erfüllt sind, so bliebe das nach diesem Grundsatz gefundene System von Werten unverändert, selbst wenn anstelle von $M$ irgendein anderer Wert $N$ beobachtet worden wäre, vorausgesetzt, dass bei der Bezeichnung $n$ für den berechneten Wert die Differenzen $M-n$ und $N-n$ dasselbe Vorzeichen haben. Der berühmte \textsc{Laplace} schränkt diesen Grundsatz allerdings durch eine zusätzliche Bedingung etwas ein. Er verlangt, dass die Summe der Differenzen selbst bei unveränderten Vorzeichen $=0$ wird. Daraus folgt, dass die Zahl der genau erfüllten Gleichungen um eins kleiner wird als die Zahl der Unbekannten. Unsere vorherige Beobachtung gilt dennoch weiterhin, sofern mindestens zwei Unbekannte vorliegen.
